\documentclass[times, a4paper, 12pt, onecolumn, oneside]{article} 

\usepackage[a4paper,left=2.0cm,right=2.0cm,bottom=1.0cm,top=2.0cm,includefoot,includehead]{geometry}
\usepackage{amsmath}
\usepackage{amsthm}
\usepackage{amsfonts}
\usepackage{amssymb}
\usepackage{graphicx} 
\usepackage{fancyhdr}
\usepackage{natbib}
\usepackage{enumitem}
\setlist{nosep}

\pagestyle{fancyplain}
\date{}
\lhead{{\em Immonen (LUT) }}
\chead{Exercises Week 3}
\rhead{\today}

\bibliographystyle{apalike}

\title{ Exercises Week 3}
\author{Veikka Immonen}
\date{\today}

\begin{document}

\maketitle

\section{Continuation of GP regression from Ex. 1}

\subsection{Generate the data from the underlying true process}

Data is from a model
$$
y = Hx + \epsilon
$$
\begin{itemize}
\item $x \sim \mathcal{GP}(0, c_\text{Matèrn}(x, x'))$
\item $c_\text{Matèrn}(x, x') = \sigma^2 \frac{2^{1 - \nu}}{\Gamma(\nu)}\left(\sqrt{2\nu}\frac{||x - x'||}{\ell}\right)^\nu K_\nu\left(\sqrt{2\nu}\frac{||x - x'||}{\ell}\right)$
\item $H$: select operator
\item $\epsilon \sim \mathcal{N}(0, \sigma_\epsilon^2)$ 
\item $n_\text{data} = 51$
\item $n_\text{state} = 201$
\end{itemize}
The true process and the sampled data presented in Figure~\ref{fig:data}, which is used throughout the Task 1.

\subsection{Hierarchical model}

For starters. Let's consider the full hierarchical model, when both the state and all hyperparameters are estimated. Let $\theta:=(\nu, \ell, \sigma, \sigma_\epsilon^2)$. The posterior of the hierarchical model is
$$
p(x, \theta | y) \propto \underbrace{p(y | x, \theta)}_\text{likelihood} \underbrace{p(x | \theta)}_\text{prior} \underbrace{p(\theta)}_\text{hyperprior}
$$
We can estimate $x$ and $\theta$ using a Metropolis-within-Gibbs (MWG) schema:
\begin{enumerate}
\item Fix $\theta$ and propose new sample $x'$ using pCN:
$$
    \alpha_\text{pCN} = \min\left\{1, \frac{p(y|x',\theta)}{p(y|x,\theta)} \right\}
$$
\item Fix $x$ and propose new sample $\theta'$ using MH:
$$
    \alpha_\text{MH} = \min\left\{1, \frac{p(y|x,\theta')p(x | \theta')p(\theta')}{p(y|x,\theta)p(x | \theta)p(\theta)} \right\}, 
$$
\end{enumerate}
assuming that the sampling is a zero-mean Gaussian random walk in both cases. Throughout the Task 1, following sampling arrangements are used:
\begin{itemize}
    \item $\theta_0 = (\nu_0, \ell_0, \sigma_0, \sigma_{\epsilon 0}^2) = (2.5, 0.1, 1.0, 0.1)$
    \item $x_0 = \text{Linear fit to } (t_i, y_i)$
    \item $\beta_x = 0.1$.
\end{itemize}
Chain length and $\beta_\theta$ are later specified for each sub-task. Note that for each non-sampled hyperparameter, the respective element from $\beta_\theta$ is set to zero.

\subsection{Sample only $x$ and $\sigma_\epsilon$}

Now the model is considerably simpler, since Matèrn covariance stays fixed. We end up with
$$
p(x, \sigma_\epsilon | y) \propto p(y | x, \sigma_\epsilon) p(x) p(\sigma_\epsilon),
$$
and the prior term cancels in $\alpha_\text{MH}$ computation. Note, that
$$
p(y|x',\sigma_\epsilon) \propto \exp\left(-\frac{1}{2\sigma_\epsilon^2} ||y-Hx||^2 \right).
$$
For this arrangement, 100000 samples of the state and hyperparameters are drawn for the chain.

\subsubsection{Sampling $\sigma_\epsilon$ from log-Gaussian}

Let $z := \log(\sigma_\epsilon^2) \sim \mathcal{N}(\mu_z, \sigma_z^2)$. We can sample $z$ directly and derive posterior as
$$
p(x, z | y) \propto p(y | x, z) p(x) p(z),
$$
where
$$
p(z) \propto \exp\left(-\frac{(\mu_z - z)^2}{2\sigma_z^2}\right).
$$
Now, we set $\mu_z = \log(0.1)$, $\sigma_z = 1.0$, and $\beta_z = 0.001$.
Results are presented in Figure~\ref{fig:1_1_1}, including mean posterior state with credible intervals and RMSE to the true process, and chains, histograms and expected values of the sampled hyperparameters.

\subsubsection{Sampling $\sigma_\epsilon$ from uniform distribution}

Define $z$ as $|z| = \sigma_\epsilon^2$ and $z \sim \mathcal{U}\left[-\infty, \infty\right]$. Again we sample $z$ directly, with $\beta_z = 0.00$. Since z is uniform (and covariance hyperparameters are fixed), the posterior is equivalent to the likelihood.
The results are presented similarly in Figure~\ref{fig:1_1_2}.

\subsection{Sample only $x$ and $\nu$, $\sigma_\epsilon$}

Now the posterior is
$$
p(x, \nu, \sigma_\epsilon | y) \propto p(y | x, \nu, \sigma_\epsilon) p(x | \nu) p(\nu, \sigma_\epsilon),
$$
yet $\alpha_\text{pCN}$ remains the same, but
$$
    \alpha_\text{MH} = \min\left\{1, \frac{p(y | x, \nu', \sigma_\epsilon') p(x | \nu') p(\nu', \sigma_\epsilon')}{p(y | x, \nu, \sigma_\epsilon) p(x | \nu) p(\nu, \sigma_\epsilon)} \right\},
$$
where
$$
p(x | \nu) \propto \exp\left(-\frac{1}{2} (x^\intercal K_\nu^{-1} x + \log |K_{\nu}|)\right)
$$
For this arrangement, 20000 samples of the state and hyperparameters are drawn for the chain.

\subsubsection{Sampling $\sigma_\epsilon$ and $\nu$ from log-Gaussian}

Let $z := \log((\nu, \sigma_\epsilon^2)) \sim \mathcal{N}(\mu_z, \mathrm{diag}\{\sigma_z^2\})$.
For the sampling, $\mu_z = \log((1.5, 0.1))$, $\sigma_z = (1.0, 1.0)$, and $\beta_z = (0.01, 0001)$.
Results are presented in Figure~\ref{fig:1_1_3}.

\subsubsection{Sampling $\nu$ and $\sigma_\epsilon$ from uniform distribution}

Define $z$ as $|z| = (\nu, \sigma_\epsilon^2)$ and $z \sim \mathcal{U}\left[-\infty, \infty\right]$. Only the hyperprior term calcels out in the $\alpha_\text{MH}$ computation.
For the sampling, $\beta_z = (0.01, 0001)$.
Results are presented in Figure~\ref{fig:1_1_4}.

\subsection{How does the prior affect the estimates?}

\begin{itemize}
    \item Solution with log-Gaussian priors were more accurate, in terms of RMSE
    \item Uniform prior caused to more regularization, resulting to an underfitted model.
    \item Noise estimates were not stable (no convergence) or accurate in either case.
    \item Matèrn smoothess estimates were stable, but converged to significantly greater value compared to the true one.
\end{itemize}

\section{From GPr to inverse problems}

\subsection{Discrete matrix approximation}

To discretize the model
$$
y(t) = \int K(t-t')x(t')dt' + \epsilon(t)
$$
to
$$
y = Kx + \epsilon
$$
Where $K$ is a convolution kernel with a local support. Now K is a boxcar kernel with parameter $a$.
$$
K_a(t - t') = \begin{cases}\frac{1}{2a}, & |t-t'| \le a \\ 0, & \text{otherwise}\end{cases}
$$
Normalization $\frac{1}{2a}$ makes the convolution a local average:
$$
y(t) = \frac{1}{2a}\int_{t-a}^{t+a}x(t')\text{d}t' + \epsilon(t).
$$
Define the underlying signal as
$$
x(t) = \begin{cases}
\exp\left(4 - \frac{25}{10t(5 - 10t)}\right), & t\in(0,0.5) \\
1, & t\in[0.7, 0.8] \\
-1, & t\in(0.8, 0.9] \\
0, & \text{otherwise}
\end{cases},
$$
For the discrete problem, take $t_i = i \Delta t,\ i=0,...,n$ (state discretization), and approximate the integral by quadrature
$$
y_i \approx \sum_{j=0}^{N-1}K_a(t_i - t_j)x_j \Delta t
$$
Therefore in our forward matrix is:
$$
K_{i,j} = K_a(t_i - t_j) \Delta t.
$$
For the data composition:
\begin{itemize}
\item $a = 0.1$
\item $\epsilon \sim \mathcal{N}(0, 0.1^2)$ 
\item $n_\text{data} = 51$
\item $n_\text{state} = 201$
\end{itemize}
True signal, and the data presented in Figure~\ref{fig:data2}.

\subsection{Sample only $x$ and $\sigma_\epsilon$}

The hierarchical model including the simplifications along the way remains pretty much the same. Only difference is that we are using a different linear operator
Also, the sampled Matèrn process is using now $\sigma=2.0$ to match the process scale better to underlying signal.

\subsubsection{Sampling $\sigma_\epsilon$ from log-Gaussian}

Results presented in Figure~\ref{fig:2_1_1}.

\subsubsection{Sampling $\sigma_\epsilon$ from uniform}

Results presented in Figure~\ref{fig:2_1_2}.

\subsection{Sample only $x$ and $\nu$, $\sigma_\epsilon$}

\subsubsection{Sampling $\sigma_\epsilon$ and $\nu$ from log-Gaussian}

Results presented in Figure~\ref{fig:2_1_3}.

\subsubsection{Sampling $\sigma_\epsilon$ and $\nu$ from uniform}

Results presented in Figure~\ref{fig:2_1_4}.

\subsection{Considering the effect of burn-in, chain length and thinning.}

Only the use of burn-in was considered during the exercise, as only the second half of the chain was used to measure mean states, their credible intervals, hyperparameter histograms and their expected values.

Did the designed burn-in produce more accurate results? Not necessarily. Again $\sigma_\epsilon$ seemed not to converge in most cases, where as $\nu$ converged almost instantly.


\section{Figures}

\begin{figure}[ht]
    \centering
    \includegraphics[width=0.8\textwidth]{./data_and_true.pdf}
    \caption{Task 1: Underlying process and the sampled data.}
    \label{fig:data}
\end{figure}

\begin{figure}[ht]
    \centering
    \includegraphics[width=0.8\textwidth]{./1_1_1_state.pdf}
    \includegraphics[width=0.8\textwidth]{./1_1_1_pchain.pdf}
    \includegraphics[width=0.8\textwidth]{./1_1_1_phist.pdf}
    \caption{Task 1: Results of sampling $\sigma_\epsilon$ from log-Gaussian.}
    \label{fig:1_1_1}
\end{figure}

\begin{figure}[ht]
    \centering
    \includegraphics[width=0.8\textwidth]{./1_1_2_state.pdf}
    \includegraphics[width=0.8\textwidth]{./1_1_2_pchain.pdf}
    \includegraphics[width=0.8\textwidth]{./1_1_2_phist.pdf}
    \caption{Task 1: Results of sampling $\sigma_\epsilon$ from uniform.}
    \label{fig:1_1_2}
\end{figure}

\begin{figure}[ht]
    \centering
    \includegraphics[width=0.8\textwidth]{./1_1_3_state.pdf}
    \includegraphics[width=0.8\textwidth]{./1_1_3_pchain.pdf}
    \includegraphics[width=0.8\textwidth]{./1_1_3_phist.pdf}
    \caption{Task 1: Results of sampling $\nu$ and $\sigma_\epsilon$ from log-Gaussian.}
    \label{fig:1_1_3}
\end{figure}

\begin{figure}[ht]
    \centering
    \includegraphics[width=0.8\textwidth]{./1_1_4_state.pdf}
    \includegraphics[width=0.8\textwidth]{./1_1_4_pchain.pdf}
    \includegraphics[width=0.8\textwidth]{./1_1_4_phist.pdf}
    \caption{Task 1: Results of sampling $\nu$ and $\sigma_\epsilon$ from uniform.}
    \label{fig:1_1_4}
\end{figure}

\begin{figure}[ht]
    \centering
    \includegraphics[width=0.8\textwidth]{./data_and_true_2.pdf}
    \caption{Task 2: Underlying signal and the sampled data.}
    \label{fig:data2}
\end{figure}

\begin{figure}[ht]
    \centering
    \includegraphics[width=0.8\textwidth]{./2_1_1_state.pdf}
    \includegraphics[width=0.8\textwidth]{./2_1_1_pchain.pdf}
    \includegraphics[width=0.8\textwidth]{./2_1_1_phist.pdf}
    \caption{Task 2: Results of sampling $\sigma_\epsilon$ from log-Gaussian.}
    \label{fig:2_1_1}
\end{figure}

\begin{figure}[ht]
    \centering
    \includegraphics[width=0.8\textwidth]{./2_1_2_state.pdf}
    \includegraphics[width=0.8\textwidth]{./2_1_2_pchain.pdf}
    \includegraphics[width=0.8\textwidth]{./2_1_2_phist.pdf}
    \caption{Task 2: Results of sampling $\sigma_\epsilon$ from uniform.}
    \label{fig:2_1_2}
\end{figure}

\begin{figure}[ht]
    \centering
    \includegraphics[width=0.8\textwidth]{./2_1_3_state.pdf}
    \includegraphics[width=0.8\textwidth]{./2_1_3_pchain.pdf}
    \includegraphics[width=0.8\textwidth]{./2_1_3_phist.pdf}
    \caption{Task 2: Results of sampling $\nu$ and $\sigma_\epsilon$ from log-Gaussian.}
    \label{fig:2_1_3}
\end{figure}

\begin{figure}[ht]
    \centering
    \includegraphics[width=0.8\textwidth]{./2_1_4_state.pdf}
    \includegraphics[width=0.8\textwidth]{./2_1_4_pchain.pdf}
    \includegraphics[width=0.8\textwidth]{./2_1_4_phist.pdf}
    \caption{Task 2: Results of sampling $\nu$ and $\sigma_\epsilon$ from uniform.}
    \label{fig:2_1_4}
\end{figure}

\end{document}
